\documentclass[letterpaper]{article}
\usepackage[submission]{aaai2027}
\usepackage[hyphens]{url}
\usepackage{graphicx}
\urlstyle{rm}
\def\UrlFont{\rm}
\usepackage{natbib}
\usepackage{caption}
\frenchspacing
\usepackage{algorithm}
\usepackage{algorithmic}
\usepackage{booktabs}

\pdfinfo{
/TemplateVersion (2027.1)
}

\setcounter{secnumdepth}{2}

\title{Verifier-Backed Heuristic Portfolios for Industrial Order Scheduling}
\author{Anonymous Submission}
\affiliations{}

\begin{document}

\maketitle

\begin{abstract}
Industrial order scheduling must jointly respect due dates, product-specific
process routes, finished-goods inventory, parallel machine capacity, worker
eligibility, and day-level worker availability.  A scheduler that only reports
success offers little protection against subtle violations of these coupled
constraints.  We present a verifier-backed industrial scheduling framework
that separates candidate construction from schedule acceptance.  The engine
first deducts inventory and expands net demand into ordered operations.  A
time-bounded portfolio of dispatching and chunked-wavefront heuristics then
constructs minute-level candidate schedules, while an independent verifier
checks quantity, precedence, due dates, worker calendars, daily labor limits,
required machine sets, and machine concurrency.  On 670 real order-level
instances, the portfolio produces 576 verifier-accepted schedules and
identifies 94 instances as infeasible under model-specific capacity lower
bounds, with no unresolved or verifier-invalid outputs.  Fixed dispatching
rules leave 28--34 instances unresolved, and smaller wavefront chunks leave
2--4 unresolved.  Complexity-stratified analysis and four audited cases
further expose the acceptance boundary.  The framework does not certify
global optimality or complete real-world infeasibility; it provides an
auditable contract for deploying heuristic schedules under modeled industrial
constraints.
\end{abstract}

\section{Introduction}

Production planners rarely schedule an isolated textbook job shop.  A single
order instance may combine multiple products and due dates, product-specific
routes, finished-goods inventory, parallel copies of a machine, and workers
who are qualified for only some operations and available on only some days.
The resulting decisions are operational: every accepted task must name a
worker, reserve all required machines, occupy a minute interval, and preserve
the route and due-date semantics of the order.  Priority rules remain useful
in such settings because they make fast, interpretable decisions, but no
single dispatching rule dominates across manufacturing conditions
\cite{panwalkar1977survey,blackstone1982survey,pinedo2016scheduling}.

This setting creates two coupled challenges.  First, constructing a feasible
schedule is combinatorial: local choices about products, workers, days, and
machines alter the capacity available to all later operations.  Second,
declaring a schedule feasible is itself a nontrivial systems obligation.  A
candidate can look plausible while violating cumulative production,
precedence, worker availability, or concurrent machine capacity.  Treating a
scheduler's self-reported status as acceptance therefore conflates search with
validation.

We address these challenges with a verifier-backed industrial heuristic
scheduling engine.  The candidate generator uses a time-bounded portfolio of
product-ordering, worker-selection, and chunked-wavefront strategies.  The
acceptance layer independently reconstructs the modeled constraints from the
order and checks the emitted task plan.  This separation permits three
semantically distinct outcomes: a verifier-accepted schedule, an infeasibility
result supported by a model-specific capacity lower bound, or an unresolved
search result.  Only the first outcome authorizes execution.

Figure~\ref{fig:overview} summarizes the resulting evidence path.  Input
normalization and inventory deduction determine net production demand.  The
portfolio then searches worker and machine calendars, and the verifier gates
the candidate before it reaches the accepted branch.  Capacity lower bounds
can terminate an instance before search, but they are not presented as a
complete infeasibility procedure.

\begin{figure}[t]
\centering
\fbox{\parbox[c][1.18in][c]{0.92\columnwidth}{\centering
Figure placeholder: orders, inventory, process routes, machines, and worker
calendars $\rightarrow$ net-demand expansion $\rightarrow$ heuristic
portfolio $\rightarrow$ independent verifier $\rightarrow$ accepted,
lower-bound infeasible, or unresolved outcome.}}
\caption{Evidence flow of the verifier-backed scheduler.  Candidate
construction and acceptance are separate stages: a feasible or zero-task
optimal candidate enters the accepted branch only after the verifier returns
\texttt{ok}.  Capacity-lower-bound outcomes contain no schedule to verify.}
\label{fig:overview}
\end{figure}

Our contributions are as follows:
\begin{itemize}
    \item We formulate an order-level industrial scheduling contract that
    combines inventory deduction, ordered production routes, parallel machine
    capacity, worker qualification, worker calendars, and due dates.
    \item We develop a time-bounded heuristic portfolio and pair it with an
    independent verifier that gates every candidate schedule against the full
    modeled constraint set.
    \item We provide an auditable 670-instance evaluation with fixed-rule and
    wavefront ablations, complexity buckets, traced case studies, and a
    separately scoped CP-SAT baseline on a stratified 50-instance cohort.
\end{itemize}

\section{Related Work}

\subsection{Dispatching Rules and Heuristic Scheduling}

Dispatching rules rank ready operations using information such as due dates,
processing times, and remaining work.  Their low decision overhead has made
them a long-standing component of reactive manufacturing systems
\cite{panwalkar1977survey,blackstone1982survey}.  Scheduling texts and surveys
also emphasize that rule effectiveness depends on the shop objective and
instance regime, motivating combinations or adaptive selection rather than a
universal rule \cite{pinedo2016scheduling}.  Our work follows this operational
line: the portfolio combines deterministic rules to improve feasible coverage
under a fixed search budget, and the experiments isolate each rule rather
than treating the portfolio as a learned policy.

\subsection{Constraint-Based Scheduling and Benchmarks}

Constraint programming represents intervals, resource capacities, and
precedence relations directly, and mature scheduling solvers support these
constructs at industrial scale \cite{laborie2018optimizer}.  CP-SAT similarly
supports discrete scheduling models and is widely used for job-shop
formulations \cite{ortools}.  Standard benchmark suites have enabled
controlled comparisons for classical scheduling problems
\cite{taillard1993benchmarks}.  The instances considered here additionally
encode finished-goods inventory, worker-day calendars, multi-machine
operations, and an executable acceptance contract.  We therefore report
CP-SAT only on a registered stratified subset and avoid equating that subset
with the full industrial evaluation.

\subsection{Verified Acceptance}

Optimization pipelines commonly separate a produced solution from the code
that evaluates feasibility and objective values.  We make this separation an
explicit system boundary.  The scheduler may choose any portfolio strategy,
but acceptance depends only on a fresh reconstruction of the order semantics
by the verifier.  This design does not turn a heuristic into an exact method;
it makes the meaning of an accepted heuristic result testable and auditable.

\section{Problem Formulation}

\subsection{Inputs and Net Demand}

An instance contains a set of dated order lines $I$, products $P$, machine
types $M$, and workers $W$.  Order line $i\in I$ requests quantity $q_i$ of
product $p(i)$ by relative day $d_i$.  Product $p$ has initial finished-goods
inventory $s_p$ and an ordered route
$R_p=(o_{p,1},\ldots,o_{p,K_p})$.  Operation $o$ has per-unit processing time
$\tau_o$, a set $M_o$ of simultaneously required machine types, and an
eligible-worker set $W_o$.  Machine type $m$ has $c_m$ interchangeable copies.
Worker $w$ is available on days $A_w$ and can work at most 480 minutes per
available day.

Inventory is consumed before production is scheduled.  For product $p$, the
net number of units requiring production is
\begin{equation}
n_p=\max\left(\sum_{i:p(i)=p}q_i-s_p,0\right).
\label{eq:net-demand}
\end{equation}
Due-date requirements are retained cumulatively after inventory deduction, so
production completed by each due day must cover the remaining demand due by
that day.  Equation~\ref{eq:net-demand} therefore removes work already covered
by stock without discarding the deadlines of the remaining demand.  When all
$n_p$ are zero, the empty schedule is trivially optimal
under the current feasibility model; this special case does not imply general
optimality.

\subsection{Schedule and Hard Constraints}

A task specifies a product and route operation, quantity, day, start and end
minute, assigned worker, and occupied machine types.  Figure~\ref{fig:task}
illustrates the semantics.  A task requiring several machine types reserves
all of them concurrently; the listed machines are not alternatives.

\begin{figure}[t]
\centering
\fbox{\parbox[c][1.02in][c]{0.92\columnwidth}{\centering
Figure placeholder: gross demand minus inventory produces net units.  Each
unit traverses operations $o_1\rightarrow o_2\rightarrow o_3$; each node
shows its worker, simultaneous machine set, and minute duration.}}
\caption{Order-to-task semantics.  Finished-goods inventory is deducted before
route expansion.  Every produced unit follows its product route in order, and
a multi-machine operation holds all required machine types over the same
interval.}
\label{fig:task}
\end{figure}

An accepted schedule must satisfy five constraint families:
\begin{enumerate}
    \item \textbf{Production and due dates.}  Completed final operations meet
    each product's net cumulative demand no later than the corresponding due
    day.
    \item \textbf{Route precedence.}  A downstream operation cannot consume
    more units than its predecessor has completed, and task duration agrees
    with operation time and quantity.
    \item \textbf{Worker qualification and availability.}  Every assigned
    worker belongs to $W_o$, is available on the task day, has no overlapping
    tasks, and accumulates at most 480 working minutes that day.
    \item \textbf{Machine requirements.}  Every task occupies exactly its
    required machine-type set $M_o$.
    \item \textbf{Machine concurrency.}  For every type $m$ and minute $t$,
    the number of overlapping tasks requiring $m$ does not exceed $c_m$.
\end{enumerate}

The primary objective in this study is accepted feasible coverage under a
time budget, not makespan or tardiness optimization.  We consequently make no
quality or global-optimality claim among multiple accepted schedules.

\section{Verifier-Backed Heuristic Portfolio}

\subsection{Prechecks and Result Semantics}

The engine first validates references among routes, machines, and workers.  It
then evaluates necessary capacity conditions over the due-date horizon.  These
conditions compare required work with total worker capacity, required
machine-type minutes with available machine capacity, and work restricted to
worker subsets with the capacity of those subsets.  If any necessary
condition fails, the engine returns \texttt{infeasible\_proven} together with
the violated lower bound.

The label is deliberately model-specific.  A lower-bound violation proves
that no schedule satisfies the encoded instance under the current capacity
semantics, but it is not a complete decision procedure and does not represent
unmodeled factory alternatives.  If all lower bounds pass and search exhausts
its budget, the result remains \texttt{no\_solution\_found} or
\texttt{time\_limit}; neither label is converted to infeasibility.

\subsection{Portfolio Candidate Construction}

The candidate generator represents each worker and machine type as a day-level
minute calendar.  For a selected unit operation, it enumerates eligible
workers and searches for the earliest interval that jointly satisfies worker
availability, daily labor capacity, and every required machine capacity.  The
chosen interval updates all affected calendars before the next operation is
placed.

The portfolio varies three decision dimensions.  Unit-ordering rules include
earliest due date, round-robin products, largest route work, and smallest route
work.  Worker rules favor the least-used worker, the most-used worker, or a
stable name order, while day search proceeds forward or backward where
applicable.  Interleaved variants advance ready route states rather than
finishing one unit at a time.  Chunked-wavefront variants group an
earliest-due sequence into chunks and repeatedly advance ready operations
across the units in each chunk.  This limits premature commitment to a single
long route while retaining deterministic calendar updates.

The timed method orders these configurations using instance scale and machine
load, then attempts them until a verifier-acceptable candidate is found or the
budget expires.  Algorithm~\ref{alg:portfolio} shows the acceptance loop.  A
strategy failure contributes diagnostic evidence but cannot authorize the
schedule.

\begin{algorithm}[t]
\caption{Verifier-backed portfolio scheduling}
\label{alg:portfolio}
\begin{algorithmic}[1]
\REQUIRE order instance $x$, time budget $B$
\STATE normalize $x$ and deduct finished-goods inventory
\IF{a necessary capacity lower bound is violated}
    \RETURN \texttt{infeasible\_proven} with the violated bound
\ENDIF
\IF{net demand is zero}
    \STATE verify the empty plan and return it only if verification succeeds
\ENDIF
\FOR{strategy $h$ in the instance-conditioned portfolio}
    \STATE construct candidate schedule $S_h$ within the remaining budget
    \IF{$S_h$ exists and $\mathrm{Verify}(x,S_h)=\texttt{ok}$}
        \RETURN accepted schedule $S_h$
    \ENDIF
\ENDFOR
\RETURN unresolved search status and diagnostics
\end{algorithmic}
\end{algorithm}

\subsection{Independent Verification}

The verifier reads the order and candidate as separate inputs; it does not
trust cached resource state or the candidate's reported status.  It validates
task fields, reconstructs inventory-adjusted production requirements, checks
route progress and due-date completion, and rebuilds worker and machine
interval usage.  A \texttt{feasible} or zero-task \texttt{optimal} candidate
counts as accepted only when verification returns \texttt{ok}.  If
verification fails, the plan is retained for diagnosis but marked invalid and
must not be executed.  Figure~\ref{fig:verification} makes this gating logic
explicit.

\begin{figure}[t]
\centering
\fbox{\parbox[c][0.98in][c]{0.92\columnwidth}{\centering
Figure placeholder: a strategy emits candidate tasks; independent checks
reconstruct quantities, precedence, workers, due dates, machine sets, and
concurrency; only the \texttt{ok} branch reaches execution.}}
\caption{Candidate-generation and verification trace.  The verifier rebuilds
constraint state from the original order and emitted tasks.  A scheduler
status alone never determines acceptance; failed checks retain structured
errors for diagnosis.}
\label{fig:verification}
\end{figure}

\section{Experiments}

\subsection{Data, Baselines, and Metrics}

We evaluate 670 real order-level JSON instances using a frozen manifest.  The
manifest contains 473 training-period, 64 development-period, and 133 test-
period instances; a 64-instance 2025-only recent group is contained within the
test period.  The scheduler does not learn from the training split.  H1--H3
therefore evaluate all 670 instances and study deterministic strategy
coverage, rather than test-set generalization of a learned model.

H1 is the portfolio timed heuristic.  H2 runs each fixed unit-ordering rule
with the common least-used-worker and forward-day choices.  H3 fixes the
chunked-wavefront mechanism and varies chunk size.  H4 runs OR-Tools CP-SAT on
a registered stratified 50-instance subset with 120 seconds per instance.  We
report H4 separately because its scope and budget differ from the full-670
heuristic experiments.

For each run, we report verifier \texttt{ok}, model-specific
\texttt{infeasible\_proven}, unresolved, verifier-invalid, and observed
elapsed time.  A schedule contributes to accepted coverage only after
verification.  All aggregate rows were recomputed from registered JSON
artifacts and passed the project's artifact QA gate.  Runtime values describe
these runs; they are not hardware-normalized algorithmic speedup claims.

\subsection{Full-670 Portfolio and Strategy Ablations}

Table~\ref{tab:full670} reports H1--H3 on the common 670-instance set.  The
portfolio yields 576 verifier-accepted schedules and 94 capacity-lower-bound
infeasibility outcomes, with no unresolved or verifier-invalid records.  In
contrast, the four fixed rules leave 28--34 instances unresolved.  The two
strongest fixed rules each accept 548 schedules, 28 fewer than the portfolio.

Chunk size also changes coverage.  Sizes 5 and 10 leave four and two instances
unresolved, respectively, whereas chunk 25 matches the portfolio's aggregate
status counts.  This result supports strategy complementarity under the
registered budget; it does not rank the quality of accepted schedules or
establish global optimality.

\begin{table*}[t]
\centering
\setlength{\tabcolsep}{5pt}
\begin{tabular}{llrrrrrr}
\toprule
Setting & H ID & Cases & OK & LB inf. & Unresolved & Invalid & Elapsed (s) \\
\midrule
Timed heuristic portfolio & H1 & 670 & \textbf{576} & 94 & \textbf{0} & 0 & 167.628 \\
Fixed earliest due & H2 & 670 & 546 & 94 & 30 & 0 & 149.063 \\
Fixed round-robin product & H2 & 670 & 548 & 94 & 28 & 0 & 139.882 \\
Fixed largest route work & H2 & 670 & 548 & 94 & 28 & 0 & 140.998 \\
Fixed smallest route work & H2 & 670 & 542 & 94 & 34 & 0 & 147.854 \\
Wavefront, chunk 5 & H3 & 670 & 572 & 94 & 4 & 0 & 157.910 \\
Wavefront, chunk 10 & H3 & 670 & 574 & 94 & 2 & 0 & 160.155 \\
Wavefront, chunk 25 & H3 & 670 & \textbf{576} & 94 & \textbf{0} & 0 & 172.358 \\
\bottomrule
\end{tabular}
\caption{Full-670 accepted-coverage evaluation for the portfolio, fixed-rule,
and chunked-wavefront settings.  \texttt{OK} denotes verifier-accepted
schedules; \texttt{LB inf.} denotes model-specific capacity-lower-bound
infeasibility.  Elapsed time is the observed run record, not a speedup claim.}
\label{tab:full670}
\end{table*}

Figure~\ref{fig:coverage} plans the same counts as a common-denominator view,
making unresolved coverage visible without introducing a schedule-quality
axis.

\begin{figure}[t]
\centering
\fbox{\parbox[c][0.90in][c]{0.92\columnwidth}{\centering
Figure placeholder: stacked bars over a common denominator of 670 show
verifier-accepted, lower-bound infeasible, and unresolved counts for H1, four
H2 rules, and three H3 chunk sizes.}}
\caption{Accepted-coverage composition for H1--H3 on the common full-670
manifest.  Fixed rules and smaller chunks change search coverage under the
registered budget; the bars do not encode schedule-quality or optimality
comparisons.}
\label{fig:coverage}
\end{figure}

\subsection{CP-SAT on the Stratified-50 Cohort}

Table~\ref{tab:cpsat120} isolates H4 from the full-670 results.  On its
stratified 50-instance cohort, CP-SAT records 44 verifier-accepted schedules
and six capacity-lower-bound infeasibility outcomes, with no unresolved or
invalid records.  These counts characterize only that subset and its
120-second per-instance limit; they are not a case-count-equivalent comparison
with Table~\ref{tab:full670}.

\begin{table}[t]
\centering
\setlength{\tabcolsep}{4pt}
\begin{tabular}{lrrrr}
\toprule
Cases & OK & LB inf. & Unresolved & Invalid \\
\midrule
50 & 44 & 6 & 0 & 0 \\
\bottomrule
\end{tabular}
\caption{CP-SAT results on the registered stratified 50-instance cohort at
120 seconds per instance.  The cohort and budget differ from the full-670
heuristic evaluation.}
\label{tab:cpsat120}
\end{table}

\subsection{Complexity and Difficulty Buckets}

H5 joins all 670 H1 cases to raw-order features and partitions each feature
into empirical tertiles.  Operation count and machine-load ratio are derived
using the scheduler's operation and capacity semantics; total work minutes
and worker-day count are directly available raw-order statistics.  The
manifest field \texttt{load\_ratio} measures worker capacity and is not used as
a substitute for machine-load ratio.

Table~\ref{tab:complexity} shows that all 94 lower-bound infeasibility outcomes
fall in the high tertiles of operation count, total work, and machine load.
Worker-day buckets distribute those outcomes more broadly.  These are
descriptive associations within H1, not causal claims about difficulty.  The
median solve time rises from 0.032 to 0.203 seconds across the operation-count
tertiles, while every bucket retains zero unresolved and zero invalid outputs.
Figure~\ref{fig:complexity} specifies the planned visual summary of these
bucketed counts and medians.

\begin{table*}[t]
\centering
\setlength{\tabcolsep}{4pt}
\begin{tabular}{llrrrr}
\toprule
Feature & Provenance & Cases & OK & LB inf. & Solve p50 (s) \\
\midrule
Operation count & Derived & 223/223/224 & 223/223/130 & 0/0/94 & 0.032/0.095/0.203 \\
Total work minutes & Available & 223/223/224 & 223/223/130 & 0/0/94 & 0.034/0.088/0.155 \\
Machine-load ratio & Derived & 223/223/224 & 223/223/130 & 0/0/94 & 0.040/0.086/0.116 \\
Worker-day count & Available & 233/214/223 & 192/182/202 & 41/32/21 & 0.065/0.077/0.097 \\
\bottomrule
\end{tabular}
\caption{Empirical-tertile analysis of the H1 full-670 run.  Each cell lists
low/medium/high buckets.  \texttt{OK} is verifier-accepted, and \texttt{LB inf.}
is model-specific capacity-lower-bound infeasibility.  The analysis is
descriptive and does not establish causality.}
\label{tab:complexity}
\end{table*}

\begin{figure}[t]
\centering
\fbox{\parbox[c][0.92in][c]{0.92\columnwidth}{\centering
Figure placeholder: four small panels show low, medium, and high tertiles for
operation count, total work, machine-load ratio, and worker-day count, with
accepted/lower-bound counts and median runtime.}}
\caption{Planned visualization of the H5 empirical-tertile analysis.  Machine
load is derived using scheduler semantics; the worker-capacity field
\texttt{load\_ratio} is not substituted for it.}
\label{fig:complexity}
\end{figure}

\subsection{Audited Verifier Cases}

Table~\ref{tab:cases} traces four H6 cases to their order, candidate solution,
and verifier records.  The first two cases stress the acceptance path with
large schedules.  The inventory case tests the zero-task boundary.  The final
case contains no schedule: \texttt{not\_applicable} means that the verifier had
nothing to check, rather than that verification passed or failed.
Figure~\ref{fig:case-study} plans a readable time window for the first complex
case; the full 7,046-operation candidate is too dense for a paper figure.

\begin{table*}[t]
\centering
\setlength{\tabcolsep}{5pt}
\begin{tabular}{p{0.15\textwidth}p{0.21\textwidth}p{0.42\textwidth}p{0.12\textwidth}}
\toprule
Category & Case & Audited evidence & Boundary \\
\midrule
Complex feasible & \texttt{SO-2025-04-0022-2} & 46,138.88 work minutes; 7,046 scheduled operations; 6,947 merged tasks. & Feasible; verifier \texttt{ok} \\
Complex feasible & \texttt{SO-2022-12-0019-2} & 9,280 scheduled operations; 9,129 merged tasks; 15.0109-second solve. & Feasible; verifier \texttt{ok} \\
Inventory / zero task & \texttt{SO-2024-10-0032-2} & Demand 5 and inventory 5; empty plan and zero verified tasks. & Current-model optimal; verifier \texttt{ok} \\
Capacity lower bound & \texttt{SO-2025-05-0003-2} & CMM CONTURA 10/16/6 requires 10,701.61 minutes but has 10,080.00 modeled capacity minutes. & \texttt{infeasible\_proven}; verifier \texttt{not\_applicable} \\
\bottomrule
\end{tabular}
\caption{Four audited cases from the H1 run.  Scheduled-operation counts come
from candidate solutions, whereas merged-task counts come from the H1 summary
or verifier record.  The lower-bound case emits no schedule.}
\label{tab:cases}
\end{table*}

\begin{figure}[t]
\centering
\fbox{\parbox[c][0.92in][c]{0.92\columnwidth}{\centering
Figure placeholder: a readable time window from
\texttt{SO-2025-04-0022-2}, plotted as worker and machine lanes with precedence
links and concurrent-capacity boundaries.}}
\caption{Planned multi-resource trace for one complex verifier-accepted case.
The figure will expose a readable window of worker and machine occupancy and
the precedence relations checked by the verifier; it will not represent the
quality distribution of all schedules.}
\label{fig:case-study}
\end{figure}

\section{Limitations and Responsible Use}

The evidence supports feasibility under the encoded constraints, not global
optimality.  Except for the inventory-covered zero-task case, the framework
provides no optimality certificate and does not compare makespan, tardiness,
or operating cost among accepted schedules.  Likewise,
\texttt{infeasible\_proven} covers necessary capacity lower bounds and selected
exact-model outcomes; it is not a complete characterization of all
infeasibility in a changing factory.

The 670 instances provide the study's principal external evidence.  Classical
job-shop benchmarks omit several modeled semantics, including finished-goods
inventory, worker-day availability, and simultaneous multi-machine
requirements, so we do not present them as interchangeable evaluations.
Actual deployment also requires data-quality controls, maintenance calendars,
rescheduling for urgent orders, and human approval.  The verifier guarantees
only the constraints it implements.

\section{Conclusion}

We presented a verifier-backed heuristic scheduling framework that turns
industrial orders into minute-level candidate plans and accepts them only
after independent constraint reconstruction.  On 670 order instances, the
portfolio covers every instance with either a verifier-accepted schedule or a
model-specific capacity-lower-bound result, while fixed-rule and chunk-size
ablations expose where strategy choice affects search coverage.

The current study does not optimize schedule quality, prove general global
optimality, or model dynamic disruptions such as breakdowns and urgent
insertions.  Extending the acceptance contract to fixed tasks and machine
outages, and adding verifier-gated repair or local search over accepted
schedules, are concrete next steps toward production deployment.

\appendix
\section{Longer CP-SAT Budget on the Registered Cohort}

H7 reruns CP-SAT on the same stratified 50-instance cohort as H4 with a
600-second per-instance limit.  It records 42 feasible schedules, two
inventory-covered zero-task optimal schedules, and six capacity-lower-bound
infeasibility outcomes.  Thus, 44 schedules are verifier \texttt{ok}, with no
unresolved or invalid outputs.  The actual p50, p95, and maximum solve times
are 2.320, 19.339, and 38.931 seconds.  Table~\ref{tab:cpsat600} is an
appendix-only run record; it is neither an equal-budget comparison with H4 nor
an equal-scope comparison with H1--H3.

\begin{table}[t]
\centering
\setlength{\tabcolsep}{4pt}
\begin{tabular}{lrrrr}
\toprule
Cases & OK & LB inf. & Unresolved & p50/p95/max (s) \\
\midrule
50 & 44 & 6 & 0 & 2.320/19.339/38.931 \\
\bottomrule
\end{tabular}
\caption{Appendix-only CP-SAT record on the registered stratified 50-instance
cohort at a 600-second per-instance limit.  Runtime cells report observed
p50/p95/maximum solve seconds.}
\label{tab:cpsat600}
\end{table}

Figure~\ref{fig:cpsat-budgets} records the planned side-by-side visualization
of H4 and H7.  Both runs cover the same cohort and produce the same aggregate
status counts; the visualization will retain the distinct nominal budgets and
observed runtime distributions rather than implying budget equivalence.

\begin{figure}[t]
\centering
\fbox{\parbox[c][0.88in][c]{0.92\columnwidth}{\centering
Figure placeholder: H4 (120-second limit) and H7 (600-second limit) on the
same stratified 50-instance cohort, showing status counts and observed runtime
quantiles with a warning against full-670 comparison.}}
\caption{Planned comparison of the two CP-SAT run records on the same
stratified cohort.  The figure will distinguish nominal limits from observed
solve times and will not compare either subset run directly with the full-670
heuristic evaluation.}
\label{fig:cpsat-budgets}
\end{figure}

\bibliography{aaai2027}

\end{document}
